\chapter{Introduction}
\label{chap:introduction}

\section{Background and Motivation}
\label{sec:motivation}
\writingnote{Introduce diabetic retinopathy and the motivation for analysing
fundus photographs. Keep the medical background relevant to the computing task.
Support medical statements with verified sources; do not reuse unsupported
statistics from example theses.}

\section{Aim and Research Questions}
\label{sec:research-questions}
\writingnote{State the aim: develop and evaluate a reproducible binary classifier
with visual attribution analysis. Formulate questions about performance under
the final evaluation protocol, the data-quality issues identified during
preparation, and what Grad-CAM and error analysis reveal or cannot establish.
Use the same questions when writing the Discussion.}

\section{Scope and Contributions}
\label{sec:contributions}
\writingnote{Describe implementation, exploratory experiments, data preparation,
final internal evaluation, and error analysis as the contributions. Distinguish
the custom CNN from established pretrained architectures. The primary task is
No DR versus DR Present; five-class grading is an exploratory extension.
Do not claim a new EfficientNet architecture or a new attribution method.}

\section{Thesis Organization}
\label{sec:organization}
\writingnote{Briefly introduce Chapters~\ref{chap:background},
\ref{chap:methods}, \ref{chap:results}, \ref{chap:discussion}, and
\ref{chap:conclusion}. Describe what each contributes rather than repeating
the table of contents.}
